\documentclass[12pt,a4paper]{article}

\usepackage{fontawesome5}
\providecommand{\faPencilAlt}{\faIcon{pencil-alt}}

\usepackage{styles/couleurs}
\usepackage{styles/mise_en_page}
\usepackage{styles/boites}
\usepackage{styles/maths}
\usepackage{styles/figures}
\usepackage{styles/evaluations}

\usepackage{enumitem}
\usepackage{array}
\usepackage{tabularx}
\usepackage{amsmath}
\usepackage{tikz}

\usetikzlibrary{calc,angles,quotes}

\setlist[itemize]{leftmargin=1.4em,itemsep=0.15em,topsep=0.2em}
\setlist[enumerate]{leftmargin=*,itemsep=0.25em,topsep=0.25em}

\begin{document}

\enteteinterro[Calculatrice autorisée]{4e}{Vocabulaire et formule du cosinus}{12 min}{10}

\exo{3}[Connaître]
Dans le triangle \(ABC\) rectangle en \(B\) : \pts{3}

\begin{enumerate}[label=\alph*)]
  \item nomme l'hypoténuse ;
  \item nomme le côté adjacent à \(\widehat A\) ;
  \item nomme le côté adjacent à \(\widehat C\).
\end{enumerate}

\lignes{3}

\exo{3}[Représenter]
Complète les relations. \pts{3}

\[
\cos(\widehat A)=\ldots\div\ldots
\]

\[
\cos(\widehat C)=\ldots\div\ldots
\]

\[
\text{côté adjacent}
=
\ldots\times\cos(\text{angle}).
\]

\lignes{3}

\exo{2}[Calculer]
Donne une valeur approchée au centième. \pts{2}

\[
\cos 35^\circ
\qquad
\cos 72^\circ.
\]

\lignes{2}

\exo{2}[Raisonner]
Explique comment vérifier rapidement que la calculatrice est en mode degré. \pts{2}

\lignes{3}

\end{document}
